\documentclass{beamer}
\usepackage[utf8]{inputenc}
\usepackage[spanish]{babel}
\usepackage{listings}
\usepackage{xcolor}
\usepackage{tikz}
\usepackage{tcolorbox}
\usepackage{fontawesome5}

% Configuración de colores - Gruvbox Light con marrón principal
\definecolor{bg0}{HTML}{fbf1c7}
\definecolor{bg1}{HTML}{ebdbb2}
\definecolor{bg2}{HTML}{d5c4a1}
\definecolor{bg3}{HTML}{bdae93}
\definecolor{fg0}{HTML}{282828}
\definecolor{fg1}{HTML}{3c3836}
\definecolor{fg2}{HTML}{504945}
\definecolor{aqua}{HTML}{689d6a}
\definecolor{aquabright}{HTML}{427b58}
\definecolor{blue}{HTML}{458588}
\definecolor{red}{HTML}{cc241d}
\definecolor{orange}{HTML}{d65d0e}
\definecolor{yellow}{HTML}{d79921}
\definecolor{green}{HTML}{98971a}
\definecolor{purple}{HTML}{b16286}
\definecolor{brown}{HTML}{AB886D}
\definecolor{brownbright}{HTML}{8B6849}

% Colores para código
\definecolor{codegreen}{HTML}{98971a}
\definecolor{codegray}{HTML}{7c6f64}
\definecolor{codepurple}{HTML}{b16286}
\definecolor{backcolour}{HTML}{f9f5d7}

% Configuración de código
\lstdefinestyle{pythonstyle}{
    backgroundcolor=\color{backcolour},   
    commentstyle=\color{codegreen},
    keywordstyle=\color{brownbright}\bfseries,
    numberstyle=\tiny\color{codegray},
    stringstyle=\color{codepurple},
    basicstyle=\ttfamily\footnotesize\color{fg0},
    breakatwhitespace=false,         
    breaklines=true,                 
    captionpos=b,                    
    keepspaces=true,                 
    numbers=left,                    
    numbersep=3pt,                  
    showspaces=false,                
    showstringspaces=false,
    showtabs=false,                  
    tabsize=2,
    language=Python,
    frame=single,
    rulecolor=\color{brown},
    xleftmargin=8pt,
    xrightmargin=8pt,
    inputencoding=utf8,
    extendedchars=true,
    literate=
        {á}{{\'a}}1 {é}{{\'e}}1 {í}{{\'i}}1 {ó}{{\'o}}1 {ú}{{\'u}}1
        {Á}{{\'A}}1 {É}{{\'E}}1 {Í}{{\'I}}1 {Ó}{{\'O}}1 {Ú}{{\'U}}1
        {ñ}{{\~n}}1 {Ñ}{{\~N}}1
        {¿}{{?`}}1 {¡}{{!`}}1
}

\lstset{style=pythonstyle}

% Tema personalizado con marrón principal
\usetheme{Madrid}
\setbeamercolor{structure}{fg=brown}
\setbeamercolor{palette primary}{bg=brown,fg=bg0}
\setbeamercolor{palette secondary}{bg=brownbright,fg=bg0}
\setbeamercolor{palette tertiary}{bg=fg1,fg=bg0}
\setbeamercolor{palette quaternary}{bg=brown,fg=bg0}
\setbeamercolor{block title}{bg=brown,fg=bg0}
\setbeamercolor{block body}{bg=bg1,fg=fg0}
\setbeamercolor{frametitle}{bg=brown,fg=bg0}
\setbeamercolor{title}{fg=fg0}
\setbeamercolor{subtitle}{fg=fg1}
\setbeamercolor{normal text}{fg=fg0,bg=bg0}
\setbeamercolor{item}{fg=brown}
\setbeamercolor{section in toc}{fg=brown}

% Cajas personalizadas con marrón
\newtcolorbox{conceptbox}[1]{
    colback=bg1,
    colframe=brown,
    fonttitle=\bfseries\small,
    title=#1,
    coltitle=fg0,
    left=3pt,
    right=3pt,
    top=3pt,
    bottom=3pt,
    boxsep=2pt
}

\newtcolorbox{notebox}{
    colback=yellow!15,
    colframe=orange,
    leftrule=2mm,
    rightrule=0mm,
    toprule=0mm,
    bottomrule=0mm,
    arc=0mm,
    left=3pt,
    right=3pt,
    top=2pt,
    bottom=2pt
}

\title{\textbf{Introducción a la Programación con Python}}
\subtitle{Semana 3: Excepciones y Manejo de Errores}
\author{$\nabla$Nablabs}
\date{}
\begin{document}

\begin{frame}
\titlepage
\end{frame}

\begin{frame}
\frametitle{Contenido}
\tableofcontents
\end{frame}

\section{¿Qué son las excepciones?}

\begin{frame}{Excepciones: Control de errores \faExclamationTriangle}
\begin{conceptbox}{¿Qué es una excepción?}
\small
\begin{itemize}
    \item[\faBug] Evento que interrumpe el flujo normal del programa
    \item[\faShieldAlt] Permite manejar errores sin que el programa se detenga
    \item[\faCode] Python levanta excepciones cuando encuentra problemas
\end{itemize}
\end{conceptbox}

\vspace{0.2cm}
\begin{center}
\large\textcolor{brown}{\faCogs\ \texttt{try}, \texttt{except}, \texttt{finally}}
\end{center}
\end{frame}

\section{Tipos de errores comunes}

\begin{frame}[fragile]{SyntaxError \faTimesCircle}
\begin{conceptbox}{Error de sintaxis}
Se produce cuando el código no sigue las reglas de Python
\end{conceptbox}

\begin{lstlisting}
# Error: falta cerrar el paréntesis
print("Hola mundo"
\end{lstlisting}

\vspace{-0.2cm}
\begin{tcolorbox}[colback=red!10,colframe=red,title={\small\faExclamationTriangle \ Error},coltitle=fg0,left=3pt,right=3pt,top=2pt,bottom=2pt,boxsep=1pt]
\small\texttt{SyntaxError: '(' was never closed}
\end{tcolorbox}

\vspace{-0.1cm}
\begin{notebox}
\small\faInfoCircle\ Los SyntaxError deben corregirse antes de ejecutar
\end{notebox}
\end{frame}

\begin{frame}[fragile]{ValueError \faExclamation}
\begin{conceptbox}{Error de valor}
Ocurre cuando una función recibe un argumento del tipo correcto pero valor inapropiado
\end{conceptbox}

\begin{lstlisting}
numero = int("abc")
\end{lstlisting}

\vspace{-0.2cm}
\begin{tcolorbox}[colback=red!10,colframe=red,title={\small\faExclamationTriangle \ Error},coltitle=fg0,left=3pt,right=3pt,top=2pt,bottom=2pt,boxsep=1pt]
\small\texttt{ValueError: invalid literal for int() with base 10: 'abc'}
\end{tcolorbox}

\vspace{-0.1cm}
\begin{notebox}
\small\faLightbulb\ Común al convertir strings a números
\end{notebox}
\end{frame}

\begin{frame}[fragile]{NameError \faQuestion}
\begin{conceptbox}{Error de nombre}
Se produce cuando se intenta usar una variable que no existe
\end{conceptbox}

\begin{lstlisting}
print(variable_no_definida)
\end{lstlisting}

\vspace{-0.2cm}
\begin{tcolorbox}[colback=red!10,colframe=red,title={\small\faExclamationTriangle \ Error},coltitle=fg0,left=3pt,right=3pt,top=2pt,bottom=2pt,boxsep=1pt]
\small\texttt{NameError: name 'variable\_no\_definida' is not defined}
\end{tcolorbox}

\vspace{-0.1cm}
\begin{notebox}
\small\faInfoCircle\ Verifica que hayas definido todas tus variables
\end{notebox}
\end{frame}

\section{Bloques try y except}

\begin{frame}[fragile]{Try-Except básico \faShieldAlt}
\begin{conceptbox}{¿Cómo funciona?}
El bloque try intenta ejecutar el código, except captura el error
\end{conceptbox}

\begin{lstlisting}
try:
    numero = int(input("Ingresa un número: "))
    print(f"El número es: {numero}")
except ValueError:
    print("Error: debes ingresar un número válido")
\end{lstlisting}

\vspace{-0.2cm}
\begin{tcolorbox}[colback=green!10,colframe=green,title={\small\faArrowCircleRight \ Salida (entrada: 'abc')},coltitle=fg0,left=3pt,right=3pt,top=2pt,bottom=2pt,boxsep=1pt]
\small\texttt{Error: debes ingresar un número válido}
\end{tcolorbox}

\vspace{-0.1cm}
\begin{notebox}
\small\faCheckCircle\ El programa continúa ejecutándose sin fallar
\end{notebox}
\end{frame}

\begin{frame}[fragile]{Múltiples excepciones \faLayerGroup}
\begin{lstlisting}
try:
    lista = [1, 2, 3]
    indice = int(input("Índice: "))
    print(lista[indice])
except ValueError:
    print("Error: debes ingresar un número")
except IndexError:
    print("Error: índice fuera de rango")
\end{lstlisting}

\vspace{-0.2cm}
\begin{columns}[T]
\column{0.48\textwidth}
\begin{tcolorbox}[colback=blue!10,colframe=blue,title={\small\faArrowDown \ Entrada: 'abc'},coltitle=fg0,left=2pt,right=2pt,top=2pt,bottom=2pt]
\small\texttt{Error: debes ingresar un número}
\end{tcolorbox}

\column{0.48\textwidth}
\begin{tcolorbox}[colback=blue!10,colframe=blue,title={\small\faArrowDown \ Entrada: 10},coltitle=fg0,left=2pt,right=2pt,top=2pt,bottom=2pt]
\small\texttt{Error: índice fuera de rango}
\end{tcolorbox}
\end{columns}
\end{frame}

\section{Cláusula else}

\begin{frame}[fragile]{Cláusula else \faCheckCircle}
\begin{conceptbox}{¿Para qué sirve?}
Se ejecuta solo si NO ocurre ninguna excepción
\end{conceptbox}

\begin{lstlisting}
try:
    numero = int(input("Ingresa un número: "))
except ValueError:
    print("Error: no es un número válido")
else:
    print(f"Perfecto! El número es: {numero}")
\end{lstlisting}

\vspace{-0.2cm}
\begin{tcolorbox}[colback=green!10,colframe=green,title={\small\faArrowCircleRight \ Salida (entrada: '42')},coltitle=fg0,left=3pt,right=3pt,top=2pt,bottom=2pt,boxsep=1pt]
\small\texttt{Perfecto! El número es: 42}
\end{tcolorbox}

\vspace{-0.1cm}
\begin{notebox}
\small\faLightbulb\ Útil para separar código que debe ejecutarse solo si no hay errores
\end{notebox}
\end{frame}

\section{Cláusula finally}

\begin{frame}[fragile]{Cláusula finally \faLock}
\begin{conceptbox}{¿Para qué sirve?}
Se ejecuta SIEMPRE, haya o no haya error
\end{conceptbox}

\begin{lstlisting}
try:
    archivo = open("datos.txt", "r")
    contenido = archivo.read()
except FileNotFoundError:
    print("Error: archivo no encontrado")
finally:
    print("Operación terminada")
    # archivo.close() si se abrió
\end{lstlisting}

\vspace{-0.2cm}
\begin{tcolorbox}[colback=green!10,colframe=green,title={\small\faArrowCircleRight \ Salida},coltitle=fg0,left=3pt,right=3pt,top=2pt,bottom=2pt,boxsep=1pt]
\small\texttt{Error: archivo no encontrado\\
Operación terminada}
\end{tcolorbox}

\vspace{-0.1cm}
\begin{notebox}
\small\faInfoCircle\ Ideal para liberar recursos (cerrar archivos, conexiones, etc.)
\end{notebox}
\end{frame}

\section{Creación de excepciones con raise}

\begin{frame}[fragile]{Lanzar excepciones con raise \faRocket}
\begin{conceptbox}{¿Cómo crear excepciones?}
Puedes lanzar excepciones manualmente con raise
\end{conceptbox}

\begin{lstlisting}
def verificar_edad(edad):
    if edad < 0:
        raise ValueError("La edad no puede ser negativa")
    if edad < 18:
        raise Exception("Debes ser mayor de edad")
    return "Acceso permitido"

try:
    print(verificar_edad(-5))
except ValueError as e:
    print(f"Error de valor: {e}")
\end{lstlisting}

\vspace{-0.2cm}
\begin{tcolorbox}[colback=red!10,colframe=red,title={\small\faArrowCircleRight \ Salida},coltitle=fg0,left=3pt,right=3pt,top=2pt,bottom=2pt,boxsep=1pt]
\small\texttt{Error de valor: La edad no puede ser negativa}
\end{tcolorbox}
\end{frame}

\section{Buenas prácticas}

\begin{frame}{Buenas prácticas en manejo de errores \faCheckDouble}
\begin{conceptbox}{Recomendaciones}
\small
\begin{itemize}
    \item[\faCheck] Captura excepciones específicas, no todas genéricas
    \item[\faCheck] No dejes bloques except vacíos
    \item[\faCheck] Usa mensajes de error descriptivos
    \item[\faCheck] Registra los errores para debugging
    \item[\faCheck] No uses excepciones para control de flujo normal
    \item[\faCheck] Limpia recursos en finally o usa context managers
\end{itemize}
\end{conceptbox}
\end{frame}

\begin{frame}[fragile]{Ejemplo: ¿Qué evitar? \faTimesCircle}
\begin{columns}[T]
\column{0.48\textwidth}
\textcolor{red}{\faTimesCircle\ Malo:}
\begin{lstlisting}
try:
    # código
    pass
except:
    pass
\end{lstlisting}

\column{0.48\textwidth}
\textcolor{green}{\faCheckCircle\ Bueno:}
\begin{lstlisting}
try:
    # código
    pass
except ValueError as e:
    print(f"Error: {e}")
\end{lstlisting}
\end{columns}

\vspace{0.3cm}
\begin{notebox}
\small\faExclamationTriangle\ Capturar todas las excepciones oculta errores importantes
\end{notebox}
\end{frame}

\section{Proyecto: Validador robusto}

\begin{frame}[fragile]{Proyecto: Validador de datos de usuario \faUserCog}
\begin{lstlisting}
def validar_usuario():
    while True:
        try:
            nombre = input("Nombre: ")
            if not nombre.strip():
                raise ValueError("El nombre no puede estar vacío")
            
            edad = int(input("Edad: "))
            if edad < 0 or edad > 120:
                raise ValueError("Edad debe estar entre 0 y 120")
            
            return {"nombre": nombre, "edad": edad}
        
        except ValueError as e:
            print(f"Error: {e}. Intenta de nuevo.")
        except KeyboardInterrupt:
            print("\nOperación cancelada")
            return None
\end{lstlisting}
\end{frame}

\begin{frame}[fragile]{Proyecto: Validador (continuación) \faUserCog}
\begin{lstlisting}
try:
    datos = validar_usuario()
    if datos:
        print(f"\nUsuario registrado:")
        print(f"  Nombre: {datos['nombre']}")
        print(f"  Edad: {datos['edad']}")
except Exception as e:
    print(f"Error inesperado: {e}")
finally:
    print("\nGracias por usar el validador")
\end{lstlisting}

\vspace{-0.2cm}
\begin{tcolorbox}[colback=green!10,colframe=green,title={\small\faArrowCircleRight \ Ejemplo de salida},coltitle=fg0,left=3pt,right=3pt,top=2pt,bottom=2pt,boxsep=1pt]
\small\texttt{Nombre: Juan\\
Edad: 25\\
Usuario registrado:\\
  Nombre: Juan\\
  Edad: 25\\
Gracias por usar el validador}
\end{tcolorbox}
\end{frame}

\begin{frame}[fragile]{Otros errores comunes \faList}
\begin{conceptbox}{Más excepciones útiles}
\small
\begin{itemize}
    \item[\faBug] \textbf{TypeError}: Operación en tipo incorrecto
    \item[\faBug] \textbf{ZeroDivisionError}: División por cero
    \item[\faBug] \textbf{KeyError}: Clave no existe en diccionario
    \item[\faBug] \textbf{FileNotFoundError}: Archivo no encontrado
    \item[\faBug] \textbf{ImportError}: Módulo no puede importarse
\end{itemize}
\end{conceptbox}

\vspace{0.2cm}
\begin{lstlisting}
try:
    resultado = 10 / 0
except ZeroDivisionError:
    print("No se puede dividir por cero")
\end{lstlisting}
\end{frame}

\begin{frame}{Resumen \faListUl}
\begin{conceptbox}{¿Qué aprendimos hoy?}
\small
\begin{itemize}
    \item[\faCheck] Tipos de errores comunes (SyntaxError, ValueError, NameError)
    \item[\faCheck] Bloques try y except
    \item[\faCheck] Cláusula else
    \item[\faCheck] Cláusula finally
    \item[\faCheck] Lanzar excepciones con raise
    \item[\faCheck] Buenas prácticas en manejo de errores
    \item[\faCheck] Proyecto: Validador robusto de datos
\end{itemize}
\end{conceptbox}

\vspace{0.3cm}
\begin{center}
\large\textcolor{brown}{\faShieldAlt\ ¡Escribe código robusto!}
\end{center}
\end{frame}

\begin{frame}
\begin{center}
\Huge\textcolor{brown}{\faQuestionCircle}\\
\vspace{0.5cm}
\Huge ¡Preguntas!\\
\vspace{1cm}
\Large\textcolor{brownbright}{Gracias por su atención}
\end{center}
\end{frame}

\end{document}
